\par\noindent\SourceLeaf{56}%
Quae series, ducta in factorem praefixum, ex formula nota, esse

\[
\sin\varphi\cos\lambda\varphi=\frac{1}{2}\sin(\lambda+1)\varphi-\frac{1}{2}\sin(\lambda-1)\varphi,
\]

\noindent solvetur in terminos, secundum sinus multiplorum procedentes. Quodlibet membrum
praeter initiale et primum in bina sic dirimitur, quorum alterum multiplum proxime
minus, alterum vero multiplum proxime maius, quam illud, cuius cosinum continet, quod
dirimitur, sortietur. Sic v. g. membrum

\[
\nu n . \sin\varphi\,A^{(2)}\cos2\varphi
\]

\noindent praebet

\[
-\nu\frac{nA^{(2)}}{2}\sin\varphi+\nu\frac{nA^{(2)}}{2}\sin3\varphi;
\]

\noindent quorum primum, adiectum illi, quod ex initiali gignitur $\nu nA\sin\varphi$, praebet seriei
evolutae terminum initialem

\[
=\left(\nu nA-\nu\frac{nA^{(2)}}{2}\right)\sin\varphi.
\]

Generatim si quaeram coefficientem, quem in serie evoluta nanciscitur $\sin(h+1)\varphi$;
soli sunt in data termini $\nu nA^{(h)}\sin\varphi\cos h\varphi$ (eiusque evoluti pars prior

\[
+\nu\frac{nA^{(h)}}{2}\sin(h+1)\varphi)
\]

\noindent ac $\nu n.A^{(h+2)}\sin\varphi\cos(h+2)\varphi$ (huius vero evoluti pars posterior

\[
-\nu\frac{nA^{(h+2)}}{2}\sin(h+1)\varphi),
\]

\noindent ex quibus membra, requisitum continentia ipsius $\varphi$ multiplum, oriri queant, estque
adeo in hac serie ipsius $\sin(h+1)\varphi$ coefficiens completus




\[
=\nu\frac{n.A^{(h)}}{2}-\nu\frac{nA^{(h+2)}}{2}.
\]

\Needspace{8\baselineskip}
$\beta$) Ex altera vero parte habebitur

\[
(1+n\cos\varphi)\left(A^{(1)}\sin\varphi+2A^{(2)}\sin2\varphi+.\,.\,.+hA^{(h)}\sin h\varphi+.\,.\,.\,.\,.\right)
\]

\noindent ex qua forma, si explicetur, duplex series originem trahit. Altera est

\[
A^{(1)}\sin\varphi+2A^{(2)}\sin2\varphi+.\,.\,.\,.+(h+1)A^{(h+1)}\sin(h+1)\varphi+.\,.\,.\,.
\]

\noindent nullo negotio prodiens.
